% Section 1: Introduction
Spatial reasoning is central to multimodal systems that operate in physical
environments: a useful VLM must determine whether one object is to the left
of another, above it, in front of it, or facing it, and it must do so using
the image---not a language prior. A growing line of work shows that
benchmark accuracy can diverge from genuine visual grounding: models can
answer correctly while depending weakly on the image~\cite{lan2026seeing,
zafar2026beyond,zhu2026visualflip}, and consistency across conditions can
coexist with error~\cite{bhat2026consistent}. Prior work has also
introduced grounding-aware training objectives that explicitly reward
visual dependence~\cite{zafar2026coral}.

What has not been examined is whether \emph{ordinary} spatial task
fine-tuning---the standard LoRA recipe with no grounding-aware
objective---moves benchmark accuracy, correct-image dependence, and
semantic pair consistency as a jointly measured vector of capability
changes, and whether those changes replicate across training seeds and
backbone families. This is the gap we address.

We use Visual Spatial Reasoning (VSR)~\cite{liu2023vsr} as the task.
VSR contains natural-image statements about object-relative relations
(left/right, above/below, facing, containment, and others) paired with
true/false labels. Following the frozen protocol of the underlying study
(documented in the repository decision log), we evaluate every checkpoint
under a corrected battery of six conditions: \emph{normal} (the frozen
prompt and image), \emph{shuffle} (frozen derangement of the image),
\emph{relcomp} (strict complement pairs: language changes, expected label
flips), \emph{facingcomp} (facing-antonym flip-law diagnostic), and the
Tier-C pair \emph{hflip\_flip} / \emph{hflip\_invariant} (global
horizontal reflection; language fixed; left/right truth must flip,
vertical/depth truth must stay).

We train three fresh General-LoRA seeds per backbone (7B \qwen{} and 2B
\smol{}) using the identical recipe as the committed seed-0 run, then
measure, per seed:
\begin{itemize}[leftmargin=*]
  \item \textbf{\dA{} (benchmark gain):} normal accuracy after tuning minus
        zero-shot normal accuracy;
  \item \textbf{\dG{} (correct-image dependence):} the change in the
        normal-minus-shuffle accuracy gap relative to zero-shot;
  \item \textbf{\dC{} (semantic pair consistency):} relcomp/facingcomp
        flip-law compliance and its seed stability;
  \item \textbf{Tier-C transformation behavior:} three separate quantities
        under global reflection---transformed accuracy
        ($A_{\text{transform}}$), pair consistency ($C_{\text{pair}}$), and
        joint correctness (both\_correct)---following the frozen metric
        definitions of the underlying study.
\end{itemize}

Every headline number in this paper was recomputed independently from raw
prediction artifacts by a standalone audit script
(\texttt{scripts/audit\_paper2\_numbers\_independent.py}) that imports no
analyzer module; the audit passed with zero discrepancies and zero
failed claims. We report means and standard deviations across the three
fresh seeds; seeds are independent training draws, not ordered stages.

Our primary finding is that the adaptation decomposes cleanly and
reproducibly. \dA{} is positive for all six fresh seeds across both
families; \dG{} is positive for all six; semantic pair consistency stays
within 0.05 of the legacy General checkpoint for every fresh seed; and all
fresh 2B Tier-C answer-update rates exceed the 2B zero-shot rate. A
post-confirmatory extension on the contemporary \qwennew{} backbone
reproduces the primary \dA{}/\dG{} pattern and the transformed-accuracy
gain, and is explicitly labeled exploratory rather than preregistered.

We make three contributions:
\begin{itemize}[leftmargin=*]
  \item \textbf{A controlled multi-seed, cross-backbone decomposition of
        ordinary spatial fine-tuning} into benchmark gain, correct-image
        dependence, and semantic pair consistency, with all quantities
        recomputed independently and frozen before manuscript mode.
  \item \textbf{A seed-robust result set:} every component of the vector
        replicates across three fresh training seeds on two backbone
        families, closing the single-checkpoint-anecdote criticism.
  \item \textbf{A scoped contemporary-backbone extension} (\qwennew{})
        that reproduces the primary adaptation pattern, honestly labeled as
        post-confirmatory external validation.
\end{itemize}
